%% LaTeX - Article customise

\documentclass[12pt]{article}
%%% PAGE DIMENSIONS
\usepackage{geometry}                % to change the page dimensions. See geometry.pdf to learn the layout options. There are lots.
\geometry{a4paper}                   % ... or a4paper or a5paper or ...
%\geometry{margins=2in} % for example, change the margins to 2 inches all round
%\geometry{landscape} % set up the page for landscape
%\usepackage[parfill]{parskip}    % Activate to begin paragraphs with an empty line rather than an indent
% read geometry.pdf for detailed page layout information

%%% PACKAGES

\usepackage{graphicx}
\usepackage{float}
\usepackage{amssymb}
\usepackage{epstopdf}
\usepackage{color}
\DeclareGraphicsRule{.tif}{png}{.png}{`convert #1 `dirname #1`/`basename #1 .tif`.png}

\usepackage{booktabs} % for much better looking tables
\usepackage{array} % for better arrays (eg matrices) in maths
\usepackage{paralist} % very flexible & customisable lists (eg. enumerate/itemize, etc.)
\usepackage{verbatim} % adds environment for commenting out blocks of text & for better verbatim
%\usepackage{subfigure} % make it possible to include more than one captioned figure/table in a single float
% These packages are all incorporated in the memoir class to one degree or another...

\usepackage{url}
\usepackage{listings}
\usepackage{longtable}
\usepackage[english]{babel}
%\usepackage[utf8]{inputenc}

\lstset{ %
%language=SQL,                % choose the language of the code
basicstyle= \small,       % the size of the fonts that are used for the code
numbers=left,                   % where to put the line-numbers
numberstyle=\footnotesize,      % the size of the fonts that are used for the line-numbers
%stepnumber=2,                   % the step between two line-numbers. If it's 1 each line
                                % will be numbered
%numbersep=5pt,                  % how far the line-numbers are from the code
%backgroundcolor=\color{white},  % choose the background color. You must add \usepackage{color}
%showspaces=false,               % show spaces adding particular underscores
%showstringspaces=false,         % underline spaces within strings
%showtabs=false,                 % show tabs within strings adding particular underscores
frame=lines,                   % adds a frame around the code
tabsize=2,                      % sets default tabsize to 2 spaces
captionpos=b,                   % sets the caption-position to bottom
breaklines=true,                % sets automatic line breaking
%breakatwhitespace=false,        % sets if automatic breaks should only happen at whitespace
title=\lstname,                 % show the filename of files included with \lstinputlisting;
                                % also try caption instead of title
escapeinside={\%*}{*)},         % if you want to add a comment within your code
morekeywords={*,function, index}            % if you want to add more keywords to the set
}

%%% HEADERS & FOOTERS
\usepackage{fancyhdr} % This should be set AFTER setting up the page geometry
%\pagestyle{fancy} % options: empty , plain , fancy
\fancypagestyle{plain}

\textwidth = 13 cm
\textheight = 21,5 cm
\oddsidemargin = 1.5 cm
\evensidemargin = 1.5 cm
\topmargin = 0.0 cm
\headheight = 0.0 cm
\headsep = 2.0 cm
\parskip = 0.5 cm
\parindent = 0.5 cm
%Thành Trương
\lhead{\scriptsize Uppsala University\\Department of Information Technology\\Th\`anh Truong\\}\chead{}\rhead{}
\selectlanguage{english}
\lfoot{}\cfoot{\thepage}\rfoot{}

\renewcommand{\headrulewidth}{0pt} % customise the layout...
\newcommand{\points}[1]{\marginpar{#1p}}
\newcommand{\uppg}[1]{\item #1}

\title{DATABASE DESIGN II - 1DL400}
\author{Assignment 3 \\ \\ Extending database index}
\date{}                                           % Activate to display a given date or no date

\begin{document}
\maketitle


\section{Extending database index}

\subsection{Objectives}

The goal of this exercise is to give a practical experience in
extending a database system with a new index structure. You will use
\emph{Meta External Index Manager (Mexima)} that provides a
general mechanism to extend AMOS II with new index implemenations.
The exercise consists of writing Java code to implement the new index
structure\emph{KDTree} and add KDTree based indexes to Amos II. You
will alos investigate the scalability of using KDTree in queries.

\subsection{Introduction}
\subsubsection{Extending database index}
When a relation \emph{R} has a large number of tuples, a query
\emph{Q} such as
\begin{lstlisting}[caption=A general query to select tuples in a given range]
select * from R where A < 'constant'
\end{lstlisting}
requires a full scan of the entire table R to select tuples that match
the predicate \texttt{A < 'constant'}. Such a full scan is expensive
to perform if R has many rows. In order to speed up the search, one
can employ \emph{B-TREE} index on attribute A, which will find the
matching tuples fast. B-trees provide scalable range queries,
i.e. matching an interval of single values against an indexed set of
single values.

However, applications often require scalabale search based on other
indexing interval search. One example is finding those vectors of
nmubers (e.g. coordinates) in a set that are close to a given vector
of numbers. B-tree may not provide scalable search of those
coordinates in a database close to a given coordinate. There are other
kinds of indexes for this; here we are going to use an index structure
called {\em K-D-trees}. You are not implementing the KD-trees
themselves, but instead tyou use a readu complied Java jar package
implementing KD-trees.

\subsubsection{Declaring indexes}
Basically, Amos II supports three kinds of indexes Linear Hashing,
B-trees, and X-trees. A specific kind of index on a parameter of a
stored function is declared by the system function:
\begin{lstlisting}[caption=Signature of AMOSQL function create\_index]
create_index(Charstring function, Charstring argname, Charstring indextype, Charstring uniqueness);
\end{lstlisting}
By default, the system defines a unique HASH index on the first
argument of stored functions.

Example:
\begin{lstlisting}[caption=Indexes on a stored function]
/*Store samples of wine quality*/
create function winequalitysamples(Charstring sp)->
	Vector of Number wq as stored;

/*Define a B-TREE index on the second argument */
create_index("winequalitysamples", "wq", "MBTREE", "multiple");

/*Examine all defined indexes */
indexes(#'winequalitysamples');
{#[OID 1362 "P_CHARSTRING.WINEQUALITYSAMPLES->VECTOR-NUMBER"],
   0,"hash","unique"}
{#[OID 1362 "P_CHARSTRING.WINEQUALITYSAMPLES->VECTOR-NUMBER"],
   1,"mbtree","multiple"}

 /*Drop a HASH index*/
 drop_index('winequalitysamples', 'sp');

/*Cannot drop the last remaining index*/
drop_index('winequalitysamples', 'wq');
Trying to remove last index on CHARSTRING.WINEQUALITYSAMPLES->VECTOR-NUMBER

/*There is only one left*/
indexes(#'winequalitysamples');
{#[OID 1311 "P_CHARSTRING.WINEQUALITYSAMPLES->VECTOR-NUMBER"],1,"mbtree","multiple"}
\end{lstlisting}

The function \texttt{indexes(\#'winequalitysamples');} shows one HASH
index by default added by the system on the first argument (position
0) and a B-TREE index on the second argument (position 1). An index
can be removed by calling
\texttt{drop\_index(Charstring functioname, Charstring
argname);}. However, the system needs at least one index defined on a
stored function. Therefore, one can not remove the last remaining
index.

\subsubsection{Foreign functions in Amos II}
Amos II users can define functions that are implemented in a
conventional programming language such as Java, C, or Lisp. One might
want to use C to efficiently handle files, memory, cache and Java to
visualize. Such Amos II functions are called \emph{foreign functions}
because they are written in an external language. Foreign functions
allow Amos II to access specialized storage managers, main memory data
structures, computational enginew, and so on. They provide the basic
primitives to extend the DBMS.

The following example gives you an example showing how to implement a
foreign function in Java. Foreign function are documented in section
6.1 ''\emph{Foreign and multi-directional functions in Amos II User
Manual}"\cite{AmosII}. There is more documentation of Amos II Java
interfaces in sections 3.1 and 3.2 of ''\emph{Foreign function
Implementaion}"in \cite{AmosIIJava}.\\

\begin{lstlisting}[label={lst:impljavakdtree}, language=Java,caption= Implementation of foreign function kdtree\_put, escapeinside={@}{@}]
/*Import Amos II Java interface packages*/
import callin.*;
import callout.*;

public class KDTreeIndex {
 . . .
/*
PUT puts <key, val> into a KDTREE given its Id. The first PUT always constructs the KDTREE
*/
	@\label{lst:impljavakdtree_sign}@ public void jkdtree_put(CallContext cxt, Tuple tpl)
		throws AmosException, KeySizeException, KeyDuplicateException {
		/*Get the id*/
		@\label{lst:impljavakdtree_getid}@int id = tpl.getIntElem(0);
		/* convert from sequence of numbers to array of floats*/
		double [] key  = toArray(tpl.getSeqElem(1));
		@\label{lst:impljavakdtree_getoid}@Oid val =  tpl.getOidElem(2);
			. . .
		/*Set the result*/
		@\label{lst:impljavakdtree_setoid}@ tpl.setElem(3, val);
		 /*Emit*/
		cxt.emit(tpl);
	}
. . .
} //end class
\end{lstlisting}

To define a foreign function in Java, one first implements the
function in Java as in Listing \ref{lst:impljavakdtree}. Then the
signature of the foreign function is defined in AmosQL and linked to
the Java implementation as in Listing
\ref{lst:defamosqlkdput}. 

\begin{lstlisting}[caption={Definition of foreign function kdtree\_put in AMOSQL},label={lst:defamosqlkdput}]
 create function kdtree_put(Integer xtId, Object f, Object o)-> Object as foreign 'JAVA:KDTreeIndex/jkdtree_put';
\end{lstlisting}

A foreign function implementation in Java always has two arguments
\emph{cxt}, and \emph{tpl} (line \ref{lst:impljavakdtree_sign}). The
argument cxt is an instance of class CallContext. It is a data
structure managed by Amos II to pass information about the caller. The
argument tpl is an instance of class Tuple and represents both
arguments and results of the foreign function call. The order of
parameters in Tuple is the same as in its signature.

Knowing the data type and position of a parameter in a foreign
function allows us to fetch or returns its value using some Amos II
Java interface functions. For example, the first input argument xtId
is accessed in Line
\ref{lst:impljavakdtree_getid} - Listing \ref{lst:impljavakdtree}
using \texttt{tpl.getIntElem(0)} since its position is 0 and its type
is Integer. The implementation returnes the value in position 3 by
\texttt{tpl.setElem} and emits the tuple tpl by
\texttt{cxt.emit(tpl)}as in Line \ref{lst:impljavakdtree_setoid} -
Listing \ref{lst:impljavakdtree}.

The following are some important methods of Amos II Java
interfaces. The rest can be referenced in \cite{AmosIIJava}.
\begin{itemize}
\item 
\texttt{int Tuple.getIntElem(int pos)} fetches an integer element at position pos.
\item 
\texttt{void Tuple.setElem(int pos, int i)} '
sets the element at position pos to integer i.
\item \texttt{Oid Tuple.getOidElem(int pos)} 
fetches an \emph{Amos object} representing any kind of data structure
which can be stored in an Amos II database.
\item \texttt{void Tuple.setElem(int pos, Oid o)}
stores Amos object o in element pos of a tuple.
\item \texttt{void CallContext.emit(Tuple)} 
emits a result tuple of a foreign function. To return a bag of
values, emit can be called interatively.
\item \texttt{Tuple Tuple.getSeqElem(int pos)}
fetches an object of type of Vector in AmosQL from position pos of a
tuple. The vector is represented in Java as a tuple objet.
\item \texttt{void Tuple.setElem(int pos, Tuple tpl)} 
stores the tuple tpl as a vector into position pos of a tuple.
\end{itemize}

\subsubsection{Mexinma}
Meta External Index Manager (Mexima), a subcomponent of Amos II,
provides a generic mechanism to extend the system with new
indexstructures. It allows the index manager in Amos II to call  an
\emph{External Index Manager (Exima)} implemeneted as a stand-alone
program in C or Java. Exima manages some special data
structure that index data in the database. To connect an
Exima to the Amos II kernel, one needs to implement a few foreign
functions, e.g. in Java. These foreign functions
resides in a program called the
\emph{Exima Driver}.

For example, in the skeleton of assignment 3, kd.jar there is a Java
package that publishes an API to interact with the KDTree index data
structure. The main program KDTreeIndex.java contains skeletons of the
Mexima foreign function implementations to access the KDTRee package.

\begin{figure}[H]
  \centering
  % Requires \usepackage{graphicx}
  \includegraphics{Mexinma1.png}
  \caption{Mexinma architecture}\label{Mexinma-architecture}
\end{figure}

Mexinma does not keep any pointer to Exinma. Instead it holds an unique identifer per Exinma (\emph{xtId}). Thus, Exinman Driver is responsible for generating xtId value when Mexinma requests to make a new instance of Exinma. Consequently, Exinma Driver maintains a list of xtId.

\begin{figure}[H]
  % Requires \usepackage{graphicx}
  \includegraphics[scale=0.75]{Mexinma_connects_Exinmas1.png}
  \caption{Mexinma connects Exinma(s) to Amos II}\label{Mexinma-connects-exinma}
\end{figure}

The developer connects an new Exima by defining some foreign functions
in two steps. In the first step, the new index type is registered to
the Mexima using the built-in function

\begin{lstlisting}[caption={Definition of register\_exindextype},label={lst:registerindextype}]
register_exindextype(Charstring indextype, Charstring langpath, Boolean saver)-> Boolean;
\end{lstlisting}

Input:

--- indextype : Type of new $<$index type$>$

--- langpath  : The string ``FOREIGN''

--- saver : The user has the option to write a special programs to
save or restore the Exema by setting the parameter {\em saver} to
true. If {\em saver} isd false Mexima will automatically provide this
service. In the excercise {\em saver} is alsways false.

Output: Success (TRUE) or failure (FALSE).

The next step is to redefine a set of foreign functions to update and
search the index structure $<$index type$>$.
\begin{itemize}
\item \textbf{create function $<$index type$>$\_make()-$>$Integer as foreign...}

creates a new index of the fiven index type and returns its index
identifier.
\item \textbf{create function $<$index type$>$\_put(Integer xtId, Object key, Object val)-$>$ Object as foreign ...}

stores a value {\em val} associated with {\em key} into an index
structure whose identifier is xtId. The value is stored as an Amos
object (class Oid).

\item \textbf{\textbf{create function $<$index type$>$\_get(Integer xtId, Object key)-$>$ Object val as foreign...}}

retrieves the value {\em val} associated with {\em key} from an index
structure whose identifier is {\em xtId}.

\item \textbf{create function $<$index type$>$\_delete(Integer xtId, Object key)-$>$ Object as foreign...}

deletes a value {\em val} associated with {\em key} from an index
structure whose identifier is {\em xtId}.

\item \textbf{create function $<$index type$>$\_clear(Integer xtId)-$>$ Object as foreign...}

removes an index structure whose identifier is {\em xtId}.

\item \textbf{create function $<$index type$>$\_load(Charstring filename)-$>$ Boolean as foreign...}

This function is optional, and is required only if the parameter
\emph{manualsaver} of function register\_exindextype is set to
true. It loads all index data of type $<$index type$>$ from a file.

\item \textbf{create function $<$index type$>$\_save(Charstring filename)-$>$ Boolean as foreign...}

This function is optional, and is required only if the parameter
\emph{saver} of function register\_exindextype is set to
true. It saves all index data of type $<$index type$>$ to files.

\item \textbf{create function $<$index type$>$\_mapper(Integer id)-$>$ Bag of Object as foreign...}\\
iterates over the index and emits pairs of keys and values. This
function is optional as not all of index structure can iterate over
all its data items such as XTREE.
\end{itemize}

Mexinma is transactional meaning that it is possible to \texttt{commit} and \texttt{rollback}. In addition, one can save and reload the system by invoking \texttt{save 'databasename.dmp'} and \texttt{javaamos databasename.dmp} respectively.

\subsubsection{K-dimensional Tree}
A K-dimensional search tree (KDTREE) is a main memory data structure
which extends binary trees to index multidimensional data.

In this assignment, a KDTREE implementation is delivered as a Java
package: kd.jar, which has an API documentation available online at
\cite{LEV11}. You need to get through that API to understand which
functions should be invoked to work with a KDTREE.

\subsection{Preparation}

Download skeleton code for assignment 3 from Assignments page
\url{http://www.it.uu.se/edu/course/homepage/dbastekn2/vt11/dbt2-vt2011-assignments.html}. and
unzip it into your home directory. There is a readme.txt file
containing instructions on how to set up your environment variables
and run the skeleton code.

For background reading you are referred
to:

\begin{itemize}
\item Elmasri and Navathe \cite{EN10}: Chapters 20, 21 and 22.
\item Padron-McCarthy and Risch \cite{PMR05}: Chapter 16.
\item All material from the lectures on object-oriented and object-relational databases systems and query languages. See also the assignment web page of this course for additional information.
\item Amos II User Manual \cite{AmosII}: Section 7.1, 6.1.
\item Amos II Java Interfaces \cite{AmosIIJava}: Section 3.1, and 3.2.
\end{itemize}
There is also an supervised introduction to the assignment in your schedule.

%\newpage
%\pagestyle{empty}
\headsep = 0.0 cm


\section{Assignment} \label{assignment}

The exercise consists of 5 parts:
\begin{enumerate}
\item Exercise 1

Task: Register with Amos II a new index type KDTREE

Instruction: see Listing \ref{lst:registerindextype}
\item  Exercise 2

Task: Redefine Mexima foreign function for the KDTree

Instruction: see Listing \ref{lst:impljavakdtree_sign}

Example:

To redefine a putter that enters a pair (key, value), you should
redefine the foreign function {\tt kdtree\_put}:

\begin{verbatim}
create function kdtree_put(Integer xtId, Object f, Object o)
-> Object as foreign 'JAVA:KDTreeIndex/kdtree_put';
\end{verbatim}

The KDTRree implementation does not support an approach to scan
through all index keys (FULL INDEX SCAN). Therefore, you do not need
to implement {\tt kdtree\_mapper}.
\item  Exercise 3

Task: Specify the index type ``KDTREE'' for the second parameter of
the stored function
\texttt{winequalitysamples})
\item  Exercise 4

Task: Investigate execution plans and test KDTREE index by answering
the following questions:
\begin{itemize}
\item What kind of index is used in the execution plan of \texttt{allsamples}?
\item Why?
\item What kind of index is used in the execution plan of \texttt{get\_sample\_example}?
\item Why?
\end{itemize}
Instruction: Method \texttt{pc('function name')} shows the execution
plan of a given a function name.
\item  Exercise 5

Task: K-nearest neighbors search (k-NN) is a method for finding K
closest points in a data set given an input point. The index structure
BTREE is usually not efficient in performing this kind of search.

In this exercise, you need to complete the implementation of
\texttt{kdtree\_knn(Vector of Number v, Number k, Function f)} which
takes v as input point, k as number of nearest neighbors, and a stored
function f storing the data set to search in.

Executing \texttt{kdtree\_knn(:v,3,\#'winequalitysamples');} raises an
error stating that\\ ``Foregin function kdtree\_knn\_search\_fn is not
yet implemented!!''

In order to utilize KDTree with KNN search, you need to complete
implementation of \texttt{kdtree\_nearest} in
\texttt{KDTreeIndex.java} then redefine this foreign function
\begin{lstlisting}[caption={Foreign function kdtree\_knn\_search\_fn},label={lst:defamosqlknn}]
   create function kdtree_knn_search_fn(Integer xtId, Vector of Number x, Number k)-> Object as foreign 'JAVA:KDTreeIndex/kdtree_nearest';
\end{lstlisting}
Now , you are be able to run \texttt{kdtree\_knn(:v,3,\#'winequalitysamples');} again.
\end{enumerate}
%\subsection{Exercises} \label{exercises}
\section{Examination}
\subsection{Hand-in documents} \label{handin}
Each exercise requires filling some AMOSQL in file
\texttt{kdtree.osql} and Java code in \texttt{KDTreeIndex.java}.

Adding descriptive comments is advised in case your solution does not
pass a validation script. Each group send only one compressed folder
(tar, zip) containing:
\begin{itemize}
\item \texttt{kdtree.osql}
\item \texttt{kdtree.lsp}
\item \texttt{KDTreeIndex.java}
\item \texttt{group.txt} stating your group number, names and emails
\end{itemize}
\subsection{Validation} \label{validation}
-- We test your solution with a validation script.\\
-- Then we check your answers for exercise 4.\\
-- Finally, we continue to investigate AMOSQL code, and Java code.\\\\
Failure on a step means a K is given . You then have to re-submit your solution until it passes that step. G is given when everything is correct.\\

%\section*{Acknowledgements}
%Contributions to the material for this assignment has been supplied by several of the former and current colleagues at the department, in alphabetical order, including Thomas Padron-McCarthy and .....

\begin{thebibliography}{ORS123}
%% <1.9> is the widest label (number or letter combination) to be used
%% in the reference list; if necessary, please change this label
%\addcontentsline{toc}{section}{References}
%\input{cdbt-vr04-ref}

%{\footnotesize
%{\small

\bibitem[EN10]{EN10} Elmasri, R. and Navathe, S. B.: Fundamentals of Databases, 6th Edition, Addison-Wesley, 2010 (available e.g. at Akademibokhandeln).

\bibitem[PMR05]{PMR05} Padron-McCarthy, T. and Risch, T.: Databasteknik, Studentlitteratur, 2005 (available e.g. at Akademibokhandeln).

\bibitem[LEV11]{LEV11} Simon D. Levy, .: KD-Tree Implementation in Java and C\#,\url{http://home.wlu.edu/~levys/software/kd/}, 2011.

\bibitem[GIST95]{GIST95} Joseph M. Hellerstein, Jeffrey F. Naughton, Avi Pfeffer, .: Generalized Search Trees for Database Systems, Proceedings of the 21st VLDB Conference Zurich, Switzerland, 1995.
\bibitem[AmosII]{AmosII} Amos II Release 13 Version 1, \url{http://user.it.uu.se/~udbl/amos/doc/amos_users_guide.html}, January 30,  2011.
\bibitem[AmosIIJava]{AmosIIJava} Amos II Java Interfaces, in subdirectories of Amos II release version downloaded at \url{http://user.it.uu.se/~torer/download/lic.php}, January 30,  2011.

%\bibitem[M10]{M10} Mimer documentation, Version 10 and more (include User's Manual, Reference Manual and Programmer's Manual), \url{http://developer.mimer.com/documentation/index.htm}.

\end{thebibliography}


%\newpage
%\section*{Appendix A: KDTreeIndex.java \label{appa}}
%
%\lstinputlisting{../KDTreeIndex.java}
%\begin{figure}[htbp]   
%\begin{center}
%  \hspace*{-1cm}
%  \includegraphics[trim=0.55cm 1.8cm 0.5cm 2.55cm, clip=true, height=9.5cm, angle=90, scale=1.25]{er-jonsonbrothers-db2.pdf}
%  \vspace*{-5cm}
%%    \caption{Entity-relationship diagram of the existing company database:}
%%\label{fig1}
%\end{center}
%\end{figure}


\begin{center}
--------------------------------------------------------------------------------------\\
\end{center}


\end{document}  